% Supplementary material for "When can AI discover adsorption kinetics?"
% A fragment to be \input into paper2/article.tex (uses \Da, \Pe, \qs, \qsm and booktabs).
% Every number carries a "% source:" comment. "% DIFF:" marks a place where the code and
% PREREG_P2_discoverability.md differ; the text reports the code.

\section*{Supplementary Table 1: column, flow and physical parameters}

Both adsorbents were solved in the same column geometry, particle size and superficial
velocity. The step case also changes the bed porosity, particle density, heat of adsorption,
feed and temperature, as recorded in erratum (b) of amendment A1. All values below are read
from the code that built each solve, and the per-cell values from the ground-truth manifest.
% source: p2_generate.py:make_physics; solver_fd.py:AdsorptionPhysicsConfig.__init__;
% fetch_real_mof_data.py:get_mof303_physics; PREREG_P2 s.9 A1 errata (b).

\begin{table}[htbp]
\centering
\small
\textbf{Supplementary Table 1.} Parameters of the two simulated columns. ``Langmuir'' is the
default physics; ``step'' is water on a MOF-303-like framework with a Type V isotherm.
\par\medskip
\begin{tabular}{@{}llll@{}}
\toprule
Quantity & Langmuir & Step & Unit \\
\midrule
\multicolumn{4}{@{}l}{\textit{Geometry and flow}}\\
Column length $L$ & 0.10 & 0.10 & m \\
Inner diameter $D_{\mathrm{in}}$ & 0.05 & 0.05 & m \\
Particle diameter $d_p$ & 0.002 & 0.002 & m \\
Superficial velocity $v$ & 0.10 & 0.10 & m\,s$^{-1}$ \\
Interstitial velocity $u = v/\varepsilon_t$ & 0.25 & 0.2857 & m\,s$^{-1}$ \\
Residence time $L/u$ & 0.40 & 0.35 & s \\
% source: solver_fd.py:AdsorptionPhysicsConfig.__init__ (L, D_in, d_p, v);
% fetch_real_mof_data.py:get_mof303_physics (L, v, d_p reset to the same values);
% u and L/u computed from these by p2_generate.py:make_physics (python one-liner, read only).
\midrule
\multicolumn{4}{@{}l}{\textit{Bed and solid}}\\
Total porosity $\varepsilon_t$ & 0.40 & 0.35 & -- \\
Particle density $\rho_p$ & 1200 & 1100 & kg\,m$^{-3}$ \\
Solid heat capacity $C_{ps}$ & 1000 & 1000 & J\,kg$^{-1}$\,K$^{-1}$ \\
Axial thermal conductivity $k_z$ & 0.10 & 0.10 & W\,m$^{-1}$\,K$^{-1}$ \\
% source: solver_fd.py:AdsorptionPhysicsConfig.__init__; fetch_real_mof_data.py:get_mof303_physics
% (rho_p and eps_t are marked UNCITED packing values there).
\midrule
\multicolumn{4}{@{}l}{\textit{Gas and thermal boundary}}\\
Gas density $\rho_g$ & 1.2 & 1.2 & kg\,m$^{-3}$ \\
Gas heat capacity $C_{pg}$ & 1000 & 1000 & J\,kg$^{-1}$\,K$^{-1}$ \\
Wall heat transfer coefficient $h_w$ & 10 & 10 & W\,m$^{-2}$\,K$^{-1}$ \\
Feed and wall temperature $T_{\mathrm{in}} = T_w$ & 298.0 & 303.15 & K \\
% source: solver_fd.py:AdsorptionPhysicsConfig.__init__; fetch_real_mof_data.py:get_mof303_physics
% (T_in = T_w = 303.15); data/p2/manifest.json key T_in (298.0 and 303.15).
\midrule
\multicolumn{4}{@{}l}{\textit{Equilibrium}}\\
Saturation capacity $q_{\max}$ & 20.0 & 25.0 & mol\,kg$^{-1}$ \\
Heat of adsorption $\Delta H$ & $-50.0$ & $-50.5$ & kJ\,mol$^{-1}$ \\
Cooperativity exponent $n$ & 1 & 4.0 & -- \\
Henry-site fraction $f_H$ & 1 & 0.06 & -- \\
Cluster affinity $b_{C0}$ & $1.720\times10^{-9}$ & $7.422\times10^{-9}$ & m$^3$\,mol$^{-1}$ \\
Primary-site affinity $b_{H0}$ & $1.720\times10^{-9}$ & $3.711\times10^{-10}$ & m$^3$\,mol$^{-1}$ \\
$b_C$ at $T_{\mathrm{in}}$ & 1.000 & 3.735 & m$^3$\,mol$^{-1}$ \\
$b_H$ at $T_{\mathrm{in}}$ & 1.000 & 0.1868 & m$^3$\,mol$^{-1}$ \\
Step position (cluster half saturation) & -- & $P/P_0 = 0.15$ at 298.15\,K & -- \\
Feed concentration $c_{\mathrm{in}}$ & 1.0 & 0.5043 (30\,\% RH) & mol\,m$^{-3}$ \\
$\qs(c_{\mathrm{in}}, T_{\mathrm{in}})$ & 10.00 & 21.90 & mol\,kg$^{-1}$ \\
$\qs(c_{\mathrm{in}}, T_{\mathrm{in}})/q_{\max}$ & 0.500 & 0.876 & -- \\
Step midpoint loading $\theta$ (L2 only) & -- & 0.3552 & -- \\
% source: isotherm.py:q_star_np (dual-term form; pure Langmuir when n = f_H = 1);
% solver_fd.py:AdsorptionPhysicsConfig.__init__ (b0 = 1/exp(50000/(R*298)), b_H0 = b0);
% fetch_real_mof_data.py:get_mof303_physics (n 4.0, f_H 0.06, b0 = calibrate_step at
% rh_to_conc(0.15, 298.15), b_H0 = b0/20, delta_H -50500, q_max 25);
% fetch_real_mof_data.py:rh_to_conc (c_in at 0.30 RH, 303.15 K);
% q* at feed, b at T_in and theta computed by isotherm.py:q_star_np and
% p2_generate.py:step_midpoint_fraction (python one-liner, read only);
% data/p2/manifest.json keys c_in (1.0, 0.5042852611085497), q_max, theta_step (0.35522387404378003).
\midrule
\multicolumn{4}{@{}l}{\textit{Kinetics, dispersion and design}}\\
Stoichiometric time $t_{\mathrm{stoich}}$ & 7200.4 & 31\,050.4 & s \\
$\Da = k\,t_{\mathrm{stoich}}$, 9 levels & \multicolumn{2}{l}{$10^{-1}$ to $10^{3}$, half-decade spacing} & -- \\
LDF coefficient $k$ & $1.389\times10^{-5}$ to $0.1389$ & $3.221\times10^{-6}$ to $0.03221$ & s$^{-1}$ \\
Residence-time form $k L/u$ & $5.56\times10^{-6}$ to $0.0556$ & $1.13\times10^{-6}$ to $0.0113$ & -- \\
Default $D_L = 0.7 D_m + 0.5\,d_p u$, $D_m = 2.5\times10^{-5}$\,m$^2$\,s$^{-1}$ & $2.675\times10^{-4}$ & $3.032\times10^{-4}$ & m$^2$\,s$^{-1}$ \\
$D_L$ multiplier (``$\Pe$'' level) & \multicolumn{2}{l}{$\times 0.2$, $\times 1$, $\times 5$} & -- \\
$D_L$ at the three levels & $5.35\times10^{-5}$, $2.675\times10^{-4}$, $1.338\times10^{-3}$ & $6.06\times10^{-5}$, $3.032\times10^{-4}$, $1.516\times10^{-3}$ & m$^2$\,s$^{-1}$ \\
$\Pe = uL/D_L$ at the three levels & 467, 93.5, 18.7 & 471, 94.2, 18.8 & -- \\
L2 rate drop and width & -- & 0.8, 0.05 & -- \\
% source: data/p2/manifest.json keys t_stoich_s (7200.4 and 31050.41), k_LDF (min, max per iso),
% Da_res (= k L/(v/eps_t), min, max per iso), D_L (three values per iso);
% p2_generate.py DA_LEVELS = logspace(-1, 3, 9), PE_MULT = (0.2, 1.0, 5.0), L2_DROP 0.8, L2_WIDTH 0.05;
% solver_fd.py:AdsorptionPhysicsConfig.stoichiometric_time and axial_dispersion (D_m 2.5e-5).
% Pe = u L / D_L is NOT stored anywhere in the code; computed here from u, L and the manifest D_L.
% DIFF: pe_mult multiplies D_L, so the cell labelled "pe0.2" has the LARGEST Peclet number
% (D_L x 0.2, Pe x 5) and "pe5" the smallest. article.tex writes the levels as "x0.2, x1, x5 of the
% default, written Pe"; the author should check that every figure and slice label reads the
% multiplier as a D_L multiplier, not a Pe multiplier.
% DIFF (already an erratum): PREREG_P2 s.4.1 says the step case changes only the isotherm;
% get_mof303_physics also sets eps_t, rho_p, delta_H, T_in and the feed (A1 errata (b)).
\midrule
\multicolumn{4}{@{}l}{\textit{Solves}}\\
Grid cells $N_z$, stored times & 400, 400 & 400, 400 & -- \\
Horizon (until outlet $c/c_{\mathrm{in}} > 0.95$) & 4 or 8 $t_{\mathrm{stoich}}$ & 4 or 8 $t_{\mathrm{stoich}}$ & -- \\
Time integrator & \multicolumn{2}{l}{BDF, rtol $10^{-6}$, atol $10^{-10}$, upwind advection} & -- \\
Cells (L1, L2, isothermal control) & 27, 0, 9 & 27, 27, 0 & -- \\
\bottomrule
\end{tabular}
\end{table}
% source: p2_generate.py N_Z = 400, N_SNAP = 400, HORIZON = 4.0, solve_cell (horizon doubled up to
% 4 times until c[-1,-1]/c_in > 0.95), cells(); solver_fd.py:generate_breakthrough_data defaults
% (rtol 1e-6, atol 1e-10, scheme "upwind", method "BDF"); data/p2/manifest.json:
% 90 cells (27 langmuir, 54 step, 9 langmuir_iso), horizon_t_stoich in {4, 8}, breakthrough_95 true
% for all 90, refined = 0 cells (the H2c refinement levels in results/p2_h2c_refine_levels.json
% are empty because no slice had a crossing, so no refined cell was solved).

The isothermal control uses the Langmuir column at $\Pe \times 1$ with no heat of adsorption.
The feed-temperature van't Hoff factor is folded into $b_{C0}$ and $b_{H0}$ first, so the
isotherm at $T_{\mathrm{in}}$ is the same as in the Langmuir cells and the bed stays at
$T_{\mathrm{in}}$.
% source: p2_generate.py:make_physics branch "langmuir_iso"; cells() (9 Da levels, Pe x1, L1).

\section*{Supplementary Note 1: methods, settings and known-answer tests}

\subsection*{Observation models}
Every method and every profile likelihood read data only through one function. O1 gives
$c$, $q$ and $T$ at 20 equally spaced probe positions. O2 gives $c$ and $T$ at the same
positions, with $q$ hidden. O3 gives the outlet $c(t)$ and $T(t)$ only. All three take 200
times, evenly spaced over the solved horizon. The probe and time indices are rounded from
a linear spacing of the stored 400 cells and 400 times.
% source: p2_observe.py:observe (N_PROBE = 20, N_T = 200; iz, it by linspace().round(); O3 iz = last cell).

\subsection*{Noise $\sigma$ and isotherm error $\varepsilon$}
Noise perturbs the observed channels only. Each observed channel gets independent Gaussian
noise with standard deviation $\sigma$ times that channel's range over the whole
ground-truth field of the cell, not over the observed points. The levels were
$\sigma \in \{0, 0.005, 0.02, 0.05\}$.
% source: p2_observe.py:observe (rng_ch = f[ch].max() - f[ch].min()); p2_grid_o1.py SIGMAS.

The isotherm error perturbs only the equilibrium the method is given. It does not touch the
data. The method receives $\qsm(c,T) = \qs(c,T)\,[1 + \varepsilon\,g(c/c_{\mathrm{in}})]$,
where $g(x) = \sum_{j=1}^{4} a_j \sin(j\pi x + \phi_j)$ divided by its root-mean-square on
2001 points of $[0,1]$, with $a_j$ standard normal and $\phi_j$ uniform on $[0, 2\pi)$. The
argument $c/c_{\mathrm{in}}$ is clipped to $[0,1]$. The function $g$ is drawn once per
(cell, replicate) and does not depend on $\sigma$ or on the observation model. The levels
were $\varepsilon \in \{0, 0.005, 0.02, 0.05\}$.
% source: p2_observe.py:isotherm_error_fn (N_MODES = 4, xs = linspace(0, 1, 2001), clip),
% _rng(name, rep, "iso"); p2_grid_o1.py EPSS.

Every draw is seeded by the first 16 hexadecimal digits of a SHA-256 hash of its keys. The
noise keys are (cell, replicate, $\sigma$, observation model); the isotherm-error keys are
(cell, replicate). So the same inputs always give the same observation. A self-test checks
that $\sigma = 0$ returns the stored field exactly, $\varepsilon = 0$ returns the true
isotherm, the noise standard deviation is within 5\,\% of $\sigma$ times the range, the
isotherm error has root-mean-square $\varepsilon$ to $10^{-3}$, draws repeat, O2 hides $q$
and O3 is outlet only.
% source: p2_observe.py:_rng, observe (r = _rng(name, rep, sigma, obs, "noise")), selftest.

Each cell had 60 replicates for every method run at the full count. Replicates 0 to 14
choose the best method of a slice; replicates 15 to 59 estimate every boundary.
% source: p2_grid_o1.py N_REP = 60; PREREG_P2 s.9 A1.1.

\subsection*{Known-answer system}
Most method tests use one synthetic system that is never a grid cell:
$\mathrm{d}q/\mathrm{d}t = k(\qs(c) - q)$ with $k = 0.02$\,s$^{-1}$, a Langmuir
$\qs(c) = 5 \cdot 3c/(1 + 3c)$, a logistic inlet ramp
$c(t) = 1/(1 + \exp(-(t - 0.4\,t_{\mathrm{end}})/(0.05\,t_{\mathrm{end}})))$,
$t_{\mathrm{end}} = 600$\,s, constant $T = 298$\,K, and 400 samples unless a test says
otherwise. It is solved with rtol $10^{-10}$ and atol $10^{-12}$.
% source: p2/synthetic.py:ldf_series (defaults n = 400, t_end = 600, qmax = 5, b = 3).

\subsection*{Shared settings: library, derivative and selection floor}
Lib-B is $\{1, c, q, T, cq, q^2, c^2, \qsm\}$, with $\qsm$ evaluated at the observed $c$ and
$T$. Lib-A is Lib-B without $\qsm$. Under L1 the true law in Lib-B is the support
$\{\qsm, q\}$. The strong-form target $\mathrm{d}q/\mathrm{d}t$ is a Savitzky--Golay
derivative with window 21 and polynomial order 3 on the uniform time grid. Rows from all
probes are stacked, so one law is fitted for the whole column.
% source: p2/features.py:lib_b, derivative (window = 21, order = 3, SciPy savgol_filter, deriv = 1);
% p2/run_o1.py:discover_o1 (rows stacked over probes).

Every sparse solver works on columns divided by their standard deviation (a constant column
keeps scale 1). Coefficients are mapped back to the original units before scoring.
% source: p2/methods.py:_standardise.

\subsection*{M2 and M2b: sparse regression in the weak and strong forms}
The strong form (M2b) regresses the Savitzky--Golay derivative on Lib-B. The weak form (M2)
multiplies the law by test functions $\phi(t) = ((t-a)(b-t))^{8}$ on windows of
\emph{support} samples, with a stride of support$/4$ (integer division), and integrates by
parts with trapezoid weights, so the target needs no derivative of the data. Each $\phi$ is
scaled to a maximum of 1. The support is swept over $\{31, 61, 121\}$ samples. For each
support the system is fitted on the probes not held out and scored by
$\|G\theta - b\|/\|b\|$ on the held-out probes, which are every fifth probe (indices 2, 7,
12 and 17 of 20). A support is admissible only if its training system has at least 3
equations per library term. The chosen support is then refitted on all probes. Both forms
run all three solvers below.
% source: p2/weak.py:weak_system (power = 8, stride = support // 4, _trapz_weights);
% p2/run_o1.py WEAK_SUPPORTS = (31, 61, 121), HOLDOUT_EVERY = 5 (held = j % 5 == 2),
% MIN_ROWS_PER_TERM = 3, choose_support, discover_o1; p2_grid_o1.py FORMS = ("strong", "weak").
% DIFF: not in PREREG_P2 s.3: if no support is admissible the code falls back to 31 samples, and if
% no probe is held out it uses 61 (p2/run_o1.py:choose_support).
% DIFF (labels only): p2/weak.py's docstring calls the weak form "M2b" and p2/methods.py's docstring
% calls STLSQ "M2"; PREREG_P2 s.3 names the weak form M2 and the strong form M2b. Results are keyed
% by solver|form, so no result depends on the label.

\paragraph{Solvers.} Sequentially thresholded least squares (STLSQ) uses ridge parameter
$10^{-6}$ and drops a term when its standardised coefficient falls below 0.05 of the largest,
for at most 20 iterations. The ensemble (M3) fits STLSQ with the same threshold on 100
bootstrap samples of 60\,\% of the rows, drawn with replacement from a fixed seed 0, keeps a
term when it appears in at least 60\,\% of the fits, and reports its median coefficient.
Exact best-subset selection (M4) enumerates every support of 1 to 4 terms, fits each by least
squares, and picks the smallest
$\mathrm{BIC} = n \ln \max(\mathrm{RSS}/n,\ (2\times10^{-3}\,\mathrm{sd}(y))^2) + m \ln n$,
with $m$ the number of terms. M4 has no sign constraints.
% source: p2/methods.py:stlsq (threshold = 0.05, alpha = 1e-6, 20 iterations),
% ensemble (n_models = 100, frac = 0.6, replace = True, inclusion = 0.6, seed = 0),
% best_subset (max_terms = 4, floor = 2e-3); PREREG_P2 s.3 table (M3, M4) and Phase A.

\paragraph{Known-answer tests.} On the known-answer system at zero noise, M4 recovers L1 with
$|\hat k - k|/k < 0.02$, and STLSQ and the ensemble both recover L1. The weak form with M4
recovers L1 at $\sigma = 0$ and at $\sigma = 2\,\%$ with $|\hat k - k|/k < 0.05$, and at
$\sigma = 2\,\%$ its $k$ error is smaller than the strong form's. On a five-probe, 200-sample
version, every solver recovers L1 in both forms through the grid's own pipeline. The
residual floor $2\times10^{-3}$ was the smallest of $\{1, 2, 5\}\times10^{-3}$ that passed at
200 samples; the derivative's error there is $8\times10^{-4}$ of $\mathrm{sd}(\mathrm{d}q/\mathrm{d}t)$.
% source: tests/test_p2.py:test_best_subset_recovers_ldf_at_zero_noise,
% test_stlsq_and_ensemble_recover_ldf_at_zero_noise, test_weak_form_recovers_ldf_at_zero_and_under_noise,
% test_strong_form_fails_where_weak_form_succeeds_under_noise,
% test_discover_o1_recovers_law_in_both_forms_and_all_solvers; p2/methods.py:best_subset docstring;
% PREREG_P2 s.3 Phase A item (i).

\subsection*{M1b: uptake from the gas balance (O2)}
On O2, $q$ is recovered from the gas-phase balance of the solver,
$\partial q/\partial t = (D_L c_{zz} - u c_z - c_t)/S$ with $u = v/\varepsilon_t$ and
$S = (1-\varepsilon_t)\rho_p/\varepsilon_t$, then integrated in time from $q = 0$ by the
trapezoid rule. $c_t$ is the Savitzky--Golay derivative above; $c_z$ and $c_{zz}$ are centred
finite differences over the probe positions. The material and flow constants are taken as
known. The two edge probes are dropped, so the O1 pipeline then runs on the inner 18 probes.
% source: p2/massbal.py:invert_uptake (solver_fd.gas_coefficients, np.gradient, cumulative_trapezoid);
% p2/run_o1.py:discover_o2 (sel = slice(1, n - 1)).
% DIFF: dropping the two edge probes on O2 is not stated in PREREG_P2 s.3 (M1b paragraph).

\paragraph{Known-answer tests.} On a fresh solver run (default physics, $k = 0.002$\,s$^{-1}$,
$N_z = 200$, horizon 3 $t_{\mathrm{stoich}}$, 300 times), the inversion recovers $q$ to within
3\,\% of $\max q$ over the inner cells. Fed the true $q$, the O2 pipeline recovers L1 with
$|\hat k - k|/k < 0.02$ (weak form, M4). With the inverted $q$ it does not recover the law at
this $\Da$, which the pre-registration declared as a pilot fact.
% source: tests/test_p2.py:test_mass_balance_inversion_recovers_q_from_dense_concentration,
% test_discover_o2_recovers_law_from_concentration_only; PREREG_P2 s.6 item 5.

\subsection*{M5: errors-in-variables best-subset (an oracle)}
M5 runs M4's enumeration of 1 to 4 terms on Lib-B in the strong form, but fits each support
by mixed least squares and total least squares. The constant column is treated as exact and
projected out by least squares. The other columns and the target are scaled by their error
standard deviations and solved by total least squares through the SVD. The selection uses
the same floored BIC as M4. The error standard deviations are the declared ones: $\sigma$
times the observed channel range for $c$, $q$ and $T$; for $\qsm$,
$\sqrt{(\varepsilon\,\qsm)^2 + (\partial\qsm/\partial c\ \sigma_c)^2}$, with the derivative
taken by a forward difference of step $10^{-6}$ times the range of $c$; first-order
propagation for $cq$, $q^2$ and $c^2$; and for the target, the $q$-channel standard
deviation times the norm of the Savitzky--Golay derivative coefficients (window 21, order 3).
At $\sigma = 0$ the target standard deviation is set to $10^{-12}$. Because M5 is given the
true error levels, it takes no part in choosing the best method for H2a, H2b or the positive
control, and is used in H2c only.
% source: p2/eiv.py:mixed_ls_tls, eiv_best_subset (max_terms = 4, floor = 2e-3);
% p2_grid_extra.py:fit_m5 (h = 1e-6 * ptp(c), savgol_coeffs(21, 3, deriv = 1), sd_y = 1e-12 at sigma 0);
% PREREG_P2 s.3 (M5) and s.9 A2.1.
% DIFF: fit_m5 builds a per-row error SD for each column, but eiv_best_subset reduces each to its
% MEAN over rows before the TLS fit (s = mean(sd_cols[n])). So the weighting is one scalar per
% column, not per observation. PREREG_P2 s.3 says "every column weighted by its DECLARED error SD",
% which this meets only in the per-column sense.
% DIFF (minor): the channel ranges in fit_m5 are taken from the OBSERVED (noisy, probe-sampled)
% data (np.ptp(ch[k])), whereas the noise SD was set from the full ground-truth field range
% (p2_observe.py:observe).

\paragraph{Known-answer tests.} With $y = 2x + 1$ and noise of standard deviation 0.15 on
$x$ (20\,000 points), ordinary least squares is attenuated by more than 0.2 in the slope and
mixed least squares and total least squares recovers slope and intercept within 0.05. At
near-zero declared error ($10^{-6}$ of each column's standard deviation) on the known-answer
system, M5 recovers L1 and selects the same support as M4. Through the grid wrapper at zero
noise on the five-probe system, M5 recovers L1 and agrees with M4.
% source: tests/test_p2.py:test_mixed_ls_tls_removes_the_attenuation_bias_ols_suffers,
% test_eiv_best_subset_recovers_ldf_and_matches_m4_at_zero_error,
% test_grid_extra_m5_matches_m4_at_zero_noise_through_the_grid_wrapper.

ODR-BINDy and WENDy were not run on the grid. The ODR-BINDy port recovered the exact
7-term support on its authors' Lorenz example in 1 of 5 noise seeds, with a mean relative
coefficient error of 0.166 against about $2.2\times10^{-3}$ reported, and is reported as
``port, accuracy below the reference''.
% source: PREREG_P2 s.3 (M5 row), s.11 (ODR-BINDy five-seed check); results/odr_bindy_seeds.json
% summary (support_rate 0.2, rel_err_mean 0.1656, reported_rel_err 0.00224).

\subsection*{M6: SINDy-PI on the isothermal control}
M6 is given no isotherm and no temperature. It runs only on the 9 isothermal-control cells,
at every $\sigma$ and at $\varepsilon = 0$, with 60 replicates, in the strong form. The
library is $\{1, c, q, cq, q^2, c^2\}$, augmented with $\mathrm{d}q/\mathrm{d}t$,
$c\,\mathrm{d}q/\mathrm{d}t$ and $q\,\mathrm{d}q/\mathrm{d}t$. Each of the three
$\mathrm{d}q$ terms is tried as the left-hand side and fitted on the rest by M4's
enumeration with up to 5 terms. The candidate with the smallest relative residual wins, and
the relation is normalised so the $\mathrm{d}q/\mathrm{d}t$ coefficient is 1. Success
requires the support $\{\mathrm{d}q, c\,\mathrm{d}q, c, q, cq\}$ exactly, with
$c\,\mathrm{d}q > 0$, $c < 0$, $q > 0$ and $cq > 0$.
% source: p2/methods.py:sindy_pi (max_terms = 5); p2_grid_extra.py:fit_m6, run (cmd m6:
% SIGMAS x (0.0,) x N_REP, strong); p2/metric.py:success_M6.

\paragraph{Known-answer test.} On the known-answer system, M6 returns exactly the support
above, and every coefficient is within 5\,\% of its true value ($c\,\mathrm{d}q = 3$,
$c = -0.3$, $q = 0.02$, $cq = 0.06$). The errors measured on 2026-10-04 were 0.02\,\% and
0.07\,\% for $b$ and $k$. Through the grid wrapper on the five-probe system, M6 succeeds.
% source: tests/test_p2.py:test_sindy_pi_finds_implicit_langmuir_ldf_without_an_isotherm,
% test_grid_extra_m6_recovers_the_implicit_law_on_the_isothermal_system; PREREG_P2 s.10 item 10.

\subsection*{M7: symbolic regression (PySR)}
M7 runs PySR with binary operators $+$, $-$ and $\times$, no unary operators, 40 iterations,
maximum expression size 12, a fixed random state equal to the replicate index, deterministic
serial search, PySR's ``best'' model selection and a time limit of 1800\,s per search. The
features are Lib-B without the constant, since PySR fits constants itself, and the target is
the strong-form derivative. Feature names are passed as placeholders $f_0, f_1, \dots$ and
mapped back. The chosen expression is expanded by SymPy. A monomial that is a single feature
to the first power counts under that feature; any other monomial is recorded as extra
support, so a spurious term makes the fit fail. M7 runs on the subgrid of all $\Da$,
$\sigma \in \{0.5, 2\}\,\%$, $\varepsilon \in \{0, 2\}\,\%$, both isotherms, $\Pe \times 1$,
O1, with 60 replicates.
% source: p2/pysr_method.py:pysr_sr (niterations = 40, maxsize = 12, timeout_s = 1800,
% deterministic = True, parallelism = "serial", model_selection = "best"), safe_names;
% p2_grid_extra.py:run (cmd m7: SUB_SIGMAS = (0.005, 0.02), SUB_EPSS = (0.0, 0.02), N_REP, Pe x1,
% features without "1", seed = rep); PREREG_P2 s.3 (M7) and s.9 A1.1.
% DIFF: PREREG_P2 s.3 says "the prereg gives every method the same features"; M7 is given Lib-B
% minus the constant column "1" (p2_grid_extra.py and p2_pilot_m7.py). PySR can still form a
% constant, which then counts as extra support.
% NOTE: pysr_sr takes an argument procs = 2 that is never used.

\paragraph{Inclusion gate.} Before any M7 search on the grid, PySR with these settings ran on
the known-answer system with 0.5\,\% relative noise on $q$ and on $\qsm$, for 3 seeds. The
rule was to include M7 if L1 was recovered in at least 2 of 3 seeds, every search finished
within 1800\,s and peak memory stayed under 3\,GB. All three seeds returned $k(\qsm - q)$ with
$\hat k = 0.019994$, $0.019992$ and $0.019943$ (relative errors $3.1\times10^{-4}$,
$4.2\times10^{-4}$ and $2.9\times10^{-3}$), in 35, 5 and 5\,s, at a peak of 0.97\,GB. M7 was
included.
% source: results/p2_pilot_m7.json (k_true 0.02, noise 0.005, limits wall_s 1800 and peak_gb 3.0,
% rows seeds 0-2, decision "INCLUDE M7 on its declared subgrid"); p2_pilot_m7.py (noise applied as
% q * (1 + 0.005 N(0,1)) and qstar * (1 + 0.005 N(0,1))); PREREG_P2 s.11.
% NOTE: the gate's noise is relative to the local value, while the grid's noise is relative to the
% channel range (p2_observe.py). PREREG_P2 s.3 only says "0.5 % noise on q AND on the measured q*".

\subsection*{M8: Kolmogorov--Arnold network (not run on the grid)}
M8 uses pykan 0.2.8 following Liu et al.'s procedure. It fits a KAN of width $[n, 1]$
(grid 5, spline order 3) on standardised raw variables $\{c, q, T, \qsm\}$ with L-BFGS for
100 steps and sparsity penalty $10^{-3}$, snaps every edge to $\{x, x^2, 0\}$, and reads the
formula back in the original variables. It is not a reimplementation of KANDy.
% source: p2/methods.py:kan_symbolic (width = [len(names), 1], grid = 5, k = 3, steps = 100,
% lamb = 1e-3, lib = ["x", "x^2", "0"]); PREREG_P2 s.3 (M8).

\paragraph{Known-answer test and cost gate.} M8 failed its known-answer test at zero noise in
every configuration tried: a polynomial library at three sparsity penalties
($10^{-3}$, $10^{-2}$, $10^{-1}$), and raw variables on one and on three trajectories. It routed
the law through $q^2$, used $q$ alone, or returned $\{1, \qsm, \qsm^2\}$. The test is kept as an
expected failure. Its grid run was gated on one fit at the grid's size (20 probes, 200 times)
taking at most 300\,s. The measured fit took 1144\,s, which projects to 68.7\,h for the 216
subgrid fits, and returned only a constant. M8 was therefore not run on the grid.
% source: tests/test_p2.py:test_kan_symbolic_recovers_ldf_at_zero_noise (xfail, strict = True);
% results/p2_cost_gates.json (limit_s 300.0, sec_per_fit 1144.44, projected_h 68.67,
% fits_on_subgrid 216, support_on_known_answer ["1"]); p2_pilot_costs.py (N_PROBE 20, N_T 200);
% PREREG_P2 s.3 (M8) and s.10.

\subsection*{M9: slow-manifold discovery}
M9 regresses the state $q$, not its rate, on the $q$-free features $\{1, c, T, c^2, \qsm\}$ by
M4's enumeration with up to 3 terms and the same floored BIC. It runs on every
non-isothermal L1 cell, O1, every $(\sigma, \varepsilon)$, 60 replicates. Success is the
support $\{\qsm\}$ with coefficient 1 within 5\,\%.
% source: p2/methods.py:slow_manifold (max_terms = 3); p2_grid_extra.py:fit_m9, run (cmd m9);
% p2/metric.py:success_manifold (tol = 0.05); PREREG_P2 s.4.3.

\paragraph{Known-answer tests.} On the known-answer system, M9 recovers $q = \qsm$ at
$k = 100$\,s$^{-1}$ and correctly does not at $k = 0.002$\,s$^{-1}$. Through the grid wrapper
on the five-probe system at $k = 20$\,s$^{-1}$, it succeeds.
% source: tests/test_p2.py:test_slow_manifold_finds_equilibrium_when_fast_and_not_when_slow,
% test_grid_extra_m9_recovers_the_slow_manifold_near_equilibrium.

\subsection*{M1: universal differential equation (demonstration)}
M1 places a network in the uptake term of the column model. The network has two hidden layers
of 32 tanh units. Its inputs are $c/c_{\mathrm{in}}$, $q/q_{\max}$, $(T - T_{\mathrm{in}})/10$\,K
and $\qsm/q_{\max}$, and its output is multiplied by $q_{\max}/t_{\mathrm{stoich}}$, which needs
no knowledge of $k$. Weights start as $0.3$ times standard normal draws from a seed equal to
the replicate index, and biases start at zero. The network is trained through a
differentiable column ($N_z = 50$, Kvaerno5, PID step control with rtol $10^{-6}$ and atol
$10^{-9}$, initial step $10^{-3}$\,s, at most 200\,000 steps). The loss is the mean squared
error of the outlet $c(t)$ and $T(t)$, each divided by its observed range. The optimiser is
SciPy's L-BFGS-B with at most 50 iterations, 60 function evaluations and 7200\,s, keeping the
best parameters seen. If a column solve exhausts its step budget after the first evaluation,
training stops and keeps the best parameters (correction B88). The trained rate over all
solved states is read out by M4 on Lib-B and by PySR with M7's settings; success on either
counts. M1 ran once: Langmuir, $\Pe \times 1$, $\sigma = 2\,\%$, $\varepsilon = 0$, O3, the 9
$\Da$ levels, one replicate.
% source: p2_grid_extra.py:_m1_rate, fit_m1 (M1_NZ = 50, M1_MAXITER = 50, M1_BUDGET_S = 7200,
% 0.3 * jax.random.normal, PRNGKey(seed)), _train_lbfgs (maxfun = maxiter + 10), run (cmd m1);
% p2/jaxcol.py:solve (Kvaerno5, PIDController rtol 1e-6 atol 1e-9, dt0 1e-3, max_steps 200_000);
% PREREG_P2 s.9 A4.
% DIFF: A4 says the inputs are "scaled by c_in, q_max and T_in"; the code feeds (T - T_in)/10,
% a shift by T_in and a fixed 10 K scale, not a division by T_in.

\paragraph{Pilot and tests.} The JAX column matches the reference solver to better than
$5\times10^{-3}$ on the exit curve, and $N_z = 50$ differs from $N_z = 400$ by 0.55\,\%. One
gradient with respect to a $2\times32$ network took 145\,s, about 50 steps in 2\,h, so M1
could not run at the declared scale and was reduced to the demonstration above. A smoke test
on a fresh column ($k = 0.002$\,s$^{-1}$) runs two L-BFGS steps and checks that the loss is
finite and the readout returns; it does not test recovery. Three tests check the B88 rule:
training keeps the best parameters when the solver fails, and a failure on the first
evaluation is still an error. M1 has no known-answer recovery test.
% source: results/p2_pilot_m1.json (exit_change_50_vs_400 0.00552, grad_nn_s 145.2,
% epochs_in_2h_at_grad_nn 49.6); tests/test_p2.py:test_jax_column_matches_verified_solver,
% test_m1_demo_trains_and_reads_out_on_a_fresh_known_answer_column; tests/test_m1_train.py;
% PREREG_P2 s.3 (M1 feasibility RESOLVED).
% NOTE: the smoke test was added with A4 (after the freeze) and the B88 tests after the grid ran,
% so "passed before the freeze" does not apply to M1's tests.

\subsection*{Success metrics}
L1 succeeds when the selected support is exactly $\{\qsm, q\}$, the $\qsm$ coefficient is
positive, the $q$ coefficient is negative, and their ratio is $-1$ within 10\,\%. The fitted
rate constant is $\hat k = -\text{coef}(q)$. For L2, with $\qsm$ fixed, the implied rate
$\hat k(q) = -(\text{coef}(q) + \text{coef}(cq)\,c + 2\,\text{coef}(q^2)\,q)$ is averaged in 10
equal bins of $q/q_{\max}$; the state dependence is found when $\{\qsm, q\}$ are selected,
the Spearman correlation over the bins is below $-0.8$, and the top bin is at most 0.6 of the
bottom bin. A fit that raised an error counts as a failure.
% source: p2/metric.py:success_L1 (ratio_tol = 0.10), k_accuracy, success_L2 (n_bins = 10,
% rho_max = -0.8, drop_max = 0.6); p2_grid_o1.py, p2_grid_extra.py (error rows success False);
% PREREG_P2 s.4.3 and s.9 A2.4.

\subsection*{Profile-likelihood identifiability}
Identifiability is judged with the true law known and with the same measured isotherm the
methods receive. $k$ is practically identifiable when the 95\,\% profile interval, the set
where $2(\max\ell - \ell) \le 3.84$, lies within $[k/2, 2k]$. All profiles refine the same
way: after the coarse grid, a pass of 21 log-spaced points is placed between the outside
neighbours of the current interval, for at most 6 passes, stopping when at least 5 points
lie inside or the step ratio falls below 1.001. The reported ends are the outside neighbours,
so the interval is conservative.
% source: p2/identifiability.py:_profile (21 points, 6 passes, okf.size >= 5, fine[1]/fine[0] < 1.001),
% identifiable; PREREG_P2 s.4.3.
% DIFF: PREREG_P2 s.3 (Phase A item ii) says the ODE profile refines with passes of 41 points and the
% outlet/PDE profile with 21. The O1 grid profile (profile_interval_o1) calls the shared _profile, so it
% uses 21 points per pass. Only the older single-probe profile_interval, used by tests, uses 41.

\paragraph{O1.} The log-likelihood is the Gaussian sum over all 20 probes. For each probe,
$\qs(t)$ is built by PCHIP interpolation of the observed $c$ and $T$ on a grid 4 times finer
than the observations, and the LDF law is integrated by its exact update for a piecewise-linear
$\qs$, vectorised over $k$. The coarse grid is $k = $ \texttt{logspace(-7, 1, 81)}, 10 points
per decade, more than a decade beyond the smallest and largest $k$ of the design. If the
coarse likelihood peaks at either end of that grid, the profile raises an error, which counts
as not identifiable. The noise standard deviation in the likelihood is
$\max(\sigma, 10^{-4})$ times the range of the observed $q$. O1 identifiability runs on every
non-isothermal L1 cell, every $(\sigma, \varepsilon)$, 60 replicates.
% source: p2/identifiability.py:loglik_o1, _qstar_path (UPSAMPLE = 4, PchipInterpolator), _ldf_exact,
% profile_interval_o1 (edge ValueError), O1_K_GRID = np.logspace(-7, 1, 81);
% p2_grid_extra.py:run (cmd o1ident, sq = max(sigma, 1e-4) * ptp(q)); PREREG_P2 s.9 A6.
% NOTE: the likelihood SD uses the range of the OBSERVED probe q (noisy), while the noise was drawn
% with SD sigma x the full-field range (p2_observe.py). PREREG_P2 does not fix this choice.

\paragraph{O2 and O3.} Each point of the profile is a full column solve at $N_z = 100$ with
the measured isotherm, scored on the interior $c$ and $T$ probes (O2) or the outlet $c$ and
$T$ (O3). The coarse grid is $k = $ \texttt{logspace(-6, 0, 25)}. The noise standard deviations
are $\max(\sigma, 10^{-4})$ times the ground-truth range of $c$ and of $T$. Each solve is capped
at 600\,s of wall time; a capped solve scores $-\infty$, and each row records the capped $k$
values and whether one lies within one coarse step ($10^{6/24}$) of the interval. These run on
the subgrid of all $\Da$, both isotherms, $\Pe \times 1$, L1, $\sigma \in \{0.5, 2\}\,\%$,
$\varepsilon \in \{0, 2\}\,\%$, with 3 replicates (216 rows per observation model).
% source: p2/identifiability.py:_profile (default k_grid = np.logspace(-6, 0, 25)),
% profile_interval_probes, profile_interval_outlet, _loglik_probes, _loglik_outlet, _with_cap,
% _solve_scored; p2_grid_o23.py SOLVE_CAP_S = 600, COARSE_STEP = 10 ** (6 / 24), run_o3, run_o2_ident
% (n_z = 100, product((0.005, 0.02), (0.0, 0.02), range(3))); PREREG_P2 s.9 A5.

\paragraph{Known-answer tests.} The exact O1 update agrees with the ODE-solver likelihood
within $2\times10^{-3}$ relative (plus 0.2) at $k \in \{0.005, 0.02, 0.05\}$. The O1 profile over 20
noisy probes ($\sigma = 2\,\%$) contains $k = 0.02$ and is identifiable. After A6, it contains $k$
at $3.22\times10^{-6}$, $1.39\times10^{-5}$ and $0.139$\,s$^{-1}$, and it raises an error on a
grid whose edge holds the peak. The single-probe profile on a fresh non-isothermal run
($k = 2\times10^{-4}$\,s$^{-1}$, $\sigma = 0.5\,\%$) lies within 2\,\% of $k$. The outlet and probe
profiles on a fresh low-$\Da$ run ($k = 2\times10^{-4}$\,s$^{-1}$) contain $k$ with
$\mathrm{hi}/\mathrm{lo} < 2$. The solve cap scores a capped solve as $-\infty$, records it, and
leaves an uncapped solve unchanged.
% source: tests/test_p2.py:test_o1_exact_ldf_likelihood_matches_the_ode_reference,
% test_o1_profile_over_all_probes_brackets_k, test_ode_profile_with_temperature_brackets_k_on_nonisothermal_probe,
% test_outlet_profile_likelihood_brackets_k_in_kinetic_regime, test_probe_profile_likelihood_brackets_k_in_kinetic_regime,
% test_profile_interval_brackets_truth_and_is_narrow_at_low_noise; tests/test_profile_range.py (A6);
% tests/test_solve_cap.py (A5).
% NOTE: the A5 and A6 tests were written during and after the grid, not before the freeze.
% \textbf{[CHECK]} the test suite was not re-run while drafting this note (machine busy); the pass
% statements follow the assertions in the test files and the outcomes recorded in PREREG_P2.
